\subsection{The exact phase of the initial motion}
The direction in which a zero starts moving is determined by a phase with a small, but essential, correction.
Write $Z(t)=e^{i\theta(t)}\zeta(1/2+it)$, with the continuous Riemann--Siegel phase normalized by $\theta(0)=0$; these are the standard functions of \citet[\S25.10(i), equations 25.10.1--2]{DLMFZeros}.

\begin{proposition}\label{prop:zero-phase}
Let $\rho=1/2+i\gamma$, $\gamma>0$, be a simple zero, and put
$\phi=\arg\zeta(3/2+i\gamma)$. Its velocity along the quotient-sheet family satisfies
\[
 \Re\rho'(0)=-\frac{|\zeta(3/2+i\gamma)|}{Z'(\gamma)}
 \left(\gamma\cos(\theta(\gamma)+\phi)
       +\frac12\sin(\theta(\gamma)+\phi)\right).
\]
Suppose the positive zeros up to this height lie on the critical line, and this zero has index $n$, counted with multiplicity. Define
\[
 W_n=n-\frac32-\frac{\theta(\gamma)+\phi}{\pi},
 \qquad \delta_n=\arctan\frac1{2\gamma}.
\]
Then the zero moves strictly to the right if and only if
\[
 \sin(\pi W_n+\delta_n)>0.
\]
For any sequence of such simple zeros with heights tending to infinity, equidistribution of $W_n$ modulo $2$ implies that the proportion moving right tends to $1/2$.
\end{proposition}
\begin{proof}
Differentiating the identity defining $Z$ at a zero gives
$\zeta'(\rho)=-i e^{-i\theta(\gamma)}Z'(\gamma)$.
Substitute this into Proposition~\ref{prop:zero-ode} and take the real part, obtaining the first formula.
Since $Z(0)=\zeta(1/2)<0$, crossing the preceding zeros, with multiplicity, gives
$\operatorname{sign}Z'(\gamma)=(-1)^{n+1}$. By the definition of $W_n$,
\[
 \cos(\theta+\phi)=(-1)^n\sin(\pi W_n),\qquad
 \sin(\theta+\phi)=(-1)^n\cos(\pi W_n).
\]
Thus
\[
 \Re\rho'(0)=\frac{|\zeta(3/2+i\gamma)|}{|Z'(\gamma)|}
 \sqrt{\gamma^2+\tfrac14}\,
 \sin(\pi W_n+\delta_n).
\]
Every factor outside the sine is positive. Finally $\delta_n\to0$. For any $\eta>0$, adding $\delta_n/\pi$ changes membership in the half-period $(0,1)$ only for finitely many initial terms or terms within $\eta$ of its boundary modulo 2. Equidistribution makes the upper density of this latter set at most $2\eta$ for small $\eta$. Letting $\eta\to0$ proves the asserted proportion.
\end{proof}

The correction $\delta_n$ cannot be dropped from an exact criterion merely because it tends to zero: points can approach the boundary of a half-period faster than $\delta_n$ does. The usual midpoint argument function satisfies $W_n=S(\gamma_n)-\phi_n/\pi$ by the zero-counting identity, but the definition above avoids any ambiguity at the zero itself.

The resulting right-moving label is a Borel function of the phase, not a character of the additive phase circle. Indeed, for $b(x)=\mathbf1_{(0,1)}(x\bmod2)$, the values at $x=y=3/10$ and at $x+y=3/5$ are all $1$, whereas $1+1=0$ in $\mathbb Z/2\mathbb Z$. A zero horizontal velocity is a third possibility and is not strict leftward motion.

The project's unpublished notes of 21 November 2026 already distinguish fixed-frequency equidistribution from a normalized central limit theorem. The following explicit counterexample records that distinction.

\begin{proposition}\label{prop:clt-not-modulo}
A Gaussian central limit theorem after normalization by an increasing standard deviation does not imply equidistribution modulo 2, even when the variance has the predicted leading term. A bounded perturbation need not preserve equidistribution modulo 2.
\end{proposition}
\begin{proof}
Let $Z$ be a standard normal random variable, let $\sigma_N\to\infty$, and set
\[
 X_N=2\operatorname{round}(\sigma_NZ/2).
\]
The rounding error has absolute value at most 1. Hence $X_N/\sigma_N$ converges in distribution to $Z$, and Cauchy--Schwarz gives
$|\operatorname{Var}(X_N)-\sigma_N^2|\le2\sigma_N+1$.
Nevertheless $X_N$ is identically 0 modulo 2. If $U$ is independent and uniform on $[0,2)$, then $Y_N=X_N+U$ is exactly uniform modulo 2 and has the same normalized Gaussian limit. The correction $Y_N-U$, bounded in magnitude by 2, destroys that uniform distribution. These are counterexamples to the stated probabilistic implications; they assert nothing about the unknown distribution of the zeta-zero phases.
\end{proof}

Consequently, a normalized central limit theorem for $S(\gamma_n)$ and boundedness of $\phi_n$ do not supply the last hypothesis of Proposition~\ref{prop:zero-phase}. Equidistribution at the fixed frequencies of $\mathbb R/2\mathbb Z$, with the arithmetic phase included, is a separate question. The limiting fraction of right-moving zeros is left open here.
